\subsection{质量表、反应Q值与两体反冲}
记核电荷数为 $Z$、中子数为 $N$、质量数为 $A=Z+N$。核反应前后先核对总电荷与总核子数，再计算质量差。若用裸核质量 $M_N$，结合能为
\[
 B=[Zm_p+(A-Z)m_n-M_N]c^2.
\]
若给的是中性原子质量 $M_{\rm at}$，并忽略电子结合能差，宜改用氢原子质量 $m_H$：
\[
 B\simeq[Zm_H+(A-Z)m_n-M_{\rm at}]c^2.
\]
两式中的质量对象必须一致；把原子质量和裸质子质量直接混在一起会多减去 $Z$ 个电子静质量。
一般核反应的 $Q=(\sum M_i-\sum M_f)c^2$，并满足 $K_f-K_i=Q$。$Q<0$ 时固定靶入射阈能通常大于 $|Q|$，因为末态还需携带总动量；精确计算用第31章的系统不变量。

当静止母核发生两体衰变，产物动量相反且大小相等。若两产物均非相对论，有 $K_1=p^2/(2m_1)$、$K_2=p^2/(2m_2)$，从而
\[
 K_1=Q\frac{m_2}{m_1+m_2},\qquad
 K_2=Q\frac{m_1}{m_1+m_2}.
\]
较轻产物带走较多动能。若一产物为光子，须用 $p_\gamma=E_\gamma/c$，不能把光子代入 $p^2/(2m)$。母体静质量 $M_i$、子体质量 $M_f$ 的两体辐射跃迁精确给出
\[
 E_\gamma=\frac{M_i^2-M_f^2}{2M_i}c^2<Q=(M_i-M_f)c^2,
\]
差值用于子体反冲。忽略反冲是建立在 $Q\ll M_ic^2$ 上的近似。
